\documentclass[10pt,a4paper]{amsart}
\usepackage[margin=19mm]{geometry}
\usepackage{amsmath,amssymb,mathtools}
\usepackage[scr=rsfs,cal=euler]{mathalfa}
\usepackage{xfrac}
\usepackage[unicode,hidelinks,bookmarks=false]{hyperref}
\newcommand{\ie}{i.e.\ }
\newcommand{\cf}{cf.\ }
\newcommand{\resp}{respectively\ }
\newcommand{\Subsection}{\subsection}
\providecommand{\nobreakdash}{\nobreak\hbox{-}}
\setlength{\emergencystretch}{2em}
\multlinegap0pt
\usepackage{tikz}
\usetikzlibrary{cd}
\usepackage{xcolor}
\usepackage{algorithm}\usepackage{algpseudocode}
\usepackage{amssymb}
\usepackage{mathtools}
\usepackage{bm}
\newtheorem{theorem}{Theorem}
\title{A Quotient Homology Theory of Representation in Neural Networks}
\author{Kosio Beshkov}
\date{Source: arXiv:2502.01360v4 (CC BY 4.0)}
\begin{document}
\maketitle

\noindent\emph{Source and licence.} A Quotient Homology Theory of Representation in Neural Networks. Version arXiv:2502.01360v4 (CC BY 4.0), distributed by arXiv under the Creative Commons Attribution 4.0 International licence, http://creativecommons.org/licenses/by/4.0/, as recorded in the arXiv licence metadata for this exact version and shown on the abstract page https://arxiv.org/abs/2502.01360v4. Full licence notice: \texttt{licenses/arxiv-2502.01360-CC-BY-4.0.txt}.\par\smallskip

\noindent\textit{Editorial scope.} Counted unit: the article's theorem proof (source0.tex lines 1040--1043). The article's complete original proof of it is delivered with it (source0.tex lines 1045--1094, one \texttt{proof} environment of 50 lines). No mathematics is written, repaired or re-derived: the statement and the proof are the article's own lines, in the article's own notation, and no retained line is substituted or re-flowed.  No further source-local context is needed: every cross-reference of every retained line resolves inside the two delivered spans.  Nothing was omitted from any retained span, so the package carries no omission record.  No retained line contains a citation, so the wrapper carries no bibliography and no bibliography entry.  The retained lines contain the article's own diagram code, which the wrapper typesets with the standard tikz packages; no image file is delivered and the package declares no asset dependency.  Editorial formatting only, in addition: an AMS \texttt{amsart} preamble that restates the article's own declarations and macros used by the retained lines, namely the environments \texttt{theorem} on separate counters, together with the standard packages those lines need.  The retained environments print in this document as Theorem 1 in order of appearance, whereas the article prints the same items with the numbering of its own full document.  The complete native source of the article as distributed in the arXiv e-print for this exact version is bundled unchanged as \texttt{sources/172651/source0.tex}.
\begin{theorem}    
    \label{Convex homology proof}
    If $\mathcal{M}\cap G_J$ is convex $\forall G_J$, then $H_k(\Phi(\mathcal{M})) \simeq H_k(\mathcal{M}/\mathcal{O}_\Phi)$.
\end{theorem}

\begin{proof}
    
    In the corollary of the previous theorem we have already shown that $H_k(\Phi(\mathcal{M})) \simeq H_k(\mathcal{M}/\sim_\Phi)$. We also know that $\mathcal{M}/\sim_{\Phi}= \mathcal{R}_\Phi \cup \mathcal{O}_\Phi\cup S - \mathcal{R}_\Phi\cap_{pair}\mathcal{O}_\Phi$, {where $S$ contains singleton sets that do not impact the topology}, therefore our proof will be complete if we can show that quotiening out the part of the rank decomposition that does not intersect the overlap decomposition leaves the homology groups invariant. We will abuse notation by writing $X-A$ instead of $X-\bigcup A$ to denote the points in $X$ that do not fall in any of the sets in $A$.

    Since we know that homology groups are invariant under homotopy equivalence, we require that $\mathcal{M}/\sim_\Phi$ is homotopy equivalent to $\mathcal{M}/\mathcal{O}_\Phi$. This is represented by the diagram:
    \begin{equation*}
        \begin{tikzcd}
        \mathcal{M} \arrow[r, "p"] \arrow[d, "q"']
        & \mathcal{M}/\mathcal{O}_\Phi \arrow[ld,shift left = 1, "\sigma"] \\
        \mathcal{M}/\sim_\Phi \arrow[ru, shift left = 1, "\pi"] \\
    \end{tikzcd}
    \end{equation*}
    The two arrows $p$ and $q$ coming out from $\mathcal{M}$ are canonical quotient maps. To show homotopy equivalence we now prove that $\pi \circ \sigma$ and $\sigma \circ \pi$ are homotopic to $\text{id}_{\mathcal{M}/\mathcal{O}_\Phi}$ and $\text{id}_{\mathcal{M}/\sim_\Phi}$ respectively. The two quotient maps split $\mathcal{M}$ into equivalence classes that we denote as $[x]_p$ and $[x]_q$. Since both $p$ and $q$ take a quotient with respect to $\mathcal{O}_\Phi$ the equivalence classes $[x]_q$ and $[x]_p$ will match on $x \in \mathcal{M} - (\mathcal{R}_\Phi-\mathcal{R}_\Phi\cap_{pair}\mathcal{O}_\Phi)$. Therefore on this set $\pi$ and $\sigma$ map identical equivalence classes to each other. The interesting part of the proof is dealing with the remaining set that is generated purely by the rank decomposition.

    We can define the map $\sigma: \mathcal{M}/\mathcal{O}_\Phi \to \mathcal{M}/\sim_\Phi$ as another quotient canonical map such that $q = \sigma \circ p$, that collapses the points that were missed by $p$ or,
    
    \begin{equation*}
        \sigma([x]_p) = 
        \begin{cases}
            q(x) \text{ if } x \in \mathcal{M} - (\mathcal{R}_\Phi-\mathcal{R}_\Phi\cap_{pair}\mathcal{O}_\Phi), \\
            q \circ p^{-1}([x]_p) \text{ if } x \in \mathcal{R}_\Phi-\mathcal{R}_\Phi\cap_{pair}\mathcal{O}_\Phi.
        \end{cases}
    \end{equation*}
    
    Note that since $p$ is injective on points that are not in $\mathcal{O}_\Phi$, it is invertible on this subdomain. The map $\pi: \mathcal{M}/\sim_\Phi \to \mathcal{M}/\mathcal{O}_\Phi$ is tricker, but we define it piecewise as,

    \begin{equation*}
    \pi([x]_q) = 
        \begin{cases}
            p(x) \text{ if } x \in \mathcal{M} - (\mathcal{R}_\Phi-\mathcal{R}_\Phi\cap_{pair}\mathcal{O}_\Phi), \\
            [z]_p \text{ such that } \sigma([z]_p) = [x]_q \text{ for some } z \in \mathcal{R}_\Phi - \mathcal{R}_\Phi\cap_{pair}\mathcal{O}_\Phi.
        \end{cases}
    \end{equation*}

    For the second condition we need to choose some representative $[z]_p$, since there are many $\sigma([z]_p) = [x]_q$, this particular choice does not really matter as in the end $\sigma \circ \pi ([x]_q) = \sigma([z]_p) = [x]_q$ and $\pi \circ \sigma([x]_p) = \pi([x]_q) =  [z]_p$. The first composition of maps is clearly homotopic to the identity, for the second we construct the homotopy.

    For the points in $\mathcal{R}_\Phi- \mathcal{R}_\Phi\cap_{pair}\mathcal{O}_\Phi$ representative of an equivalence class $[x]_q$ we can define a strong deformation retraction as $F:\mathcal{M}/\mathcal{O}_\Phi\times [0,1] \to \mathcal{M}/\mathcal{O}_\Phi$, where we have $F(x,0) = \text{id}_{\mathcal{M}/\mathcal{O}_\Phi}([x]_p)$ and $F(x,1) = \pi \circ \sigma ([x]_p)$. 
    \begin{equation*}
        F([x]_p,t) = 
        \begin{cases}
            [x]_p \text{ if } x\notin \mathcal{R}_\Phi- \mathcal{R}_\Phi\cap_{pair}\mathcal{O}_\Phi),\\
            \gamma_z(t)([x]_p) \text{ if } x \in \mathcal{R}_\Phi- \mathcal{R}_\Phi\cap_{pair}\mathcal{O}_\Phi).
        \end{cases}
    \end{equation*}

    Where $\gamma_z(t)$ is a continuous (in both $t$ and $z$) path from any point $[x]_p \in p\circ q^{-1}([x]_q)$ to the representative $[z]_p$ point. This map starts at $\gamma_z(0)([x]_p) = [x]_p$ and ends at $\gamma_z(1)([x]_p) = [z]_p$ while following a continuous path. Such a path is guaranteed to exist due to the fact that equivalence classes in $\mathcal{R}_\Phi$ arise as a result of a low-rank projection of a convex set, which implies they are convex (and therefore contractible) in $\mathcal{M}/\mathcal{O}_\Phi$ themselves. This proves that there is a homotopy equivalence between $\mathcal{M}/\sim_\Phi$ and $\mathcal{M}/\mathcal{O}_\Phi$. Since homology groups are invariant under homotopy equivalence we get our final result,
    \begin{equation*}
        H_k(\mathcal{M}/\sim_\Phi) \simeq H_k(\mathcal{M}/\mathcal{O}_\Phi).
    \end{equation*}
\end{proof}

\end{document}
